\subsection{Arithmetic domain examples}

The following small matrices exhibit the obstruction in nonmaximal
orders of totally ramified fields.

\begin{example}\label{ex:cubic-depth-two-nine}
Take $p=3$, $d=9$, $\alpha=4$ and $c=-7$.  Then
$\alpha^3+9c=1$, and over $A=\Z_3[t]/(t^3-9)$ the matrix
\[
 T=\begin{pmatrix}
 4&0&-7t\\ t^2&0&-16\\0&1&-4
 \end{pmatrix}
\]
satisfies $T^3=I$.  The residue parameter is $2$, whose degree-one
minimal polynomial has a free scalar orbit under
$\F_3^\times$.  The corresponding coordinate shows that
$[\operatorname{Ind}_{C_3}^{S_3}L]-[A[S_3]]$ has exact order
three.  Finite unramified parameter extensions give the
countable family for this same coefficient ring.

If $\pi^3=3$, then $t=\pi^2$ has valuation $2/3$, so
$t^3-9$ is irreducible and its field is the totally ramified
cubic field $\Q_3(\pi)$.  Its order is
\[
 \Z_3[\pi^2]=\Z_3\cdot1\oplus\Z_3\cdot\pi^2
                         \oplus\Z_3\cdot3\pi,
\]
of index three in the maximal order $\Z_3[\pi]$.
\end{example}

\begin{example}\label{ex:deep-cubic-binomial}
Take $p=3$, $n=3$, and $d=81$.  With $\alpha=28$ and $c=-271$,
\[
 T=\begin{pmatrix}
 28&0&-271t\\ t^2&0&-784\\0&1&-28
 \end{pmatrix},\qquad T^3=I.
\]
Indeed $28^3+81(-271)=1$.  The induced class minus the regular
class has exact order three: modulo $3$ its parameter is $2$,
and the corresponding truncated-polynomial coordinate is one.
The coefficient ring is a totally ramified cubic local domain.
If $\pi^3=3$, then $t=3\pi$ satisfies $t^3=81$ and has valuation
$4/3$; its degree over $\Q_3$ is therefore three.  The infinite
family in the corollary is obtained by the indicated finite
coefficient extensions.
\end{example}

\begin{example}\label{ex:quartic-binomial-order}
Let $O=\Z_{(3)}$ or $\Z_3$, and put
\[
 A=O[t]/(t^4-27).
\]
The choices $\alpha=-8$ and $\kappa=1/19$ give
\[
 T=\begin{pmatrix}
 -8&0&19t\\
 t^3&0&-64\\
 0&1&8
 \end{pmatrix},\qquad T^3=I.
\]
Indeed $\alpha^3=1-27\cdot19$, and its characteristic polynomial
is $z^3+512-19t^4=z^3-1$.  Restriction to
$O[s]/(s^2-27)$, $s=t^2$, has the nonzero coordinate $\mu_{u-1}$,
since $19\equiv1\pmod3$.  The induced degree-three class is nonzero,
has exact order four and vanishes generically.  Indeed, the
stable-module coordinate in Theorem~\ref{thm:binomial-depth-two-all-exponents}
is one modulo four: its residue parameter is $1$, of degree one,
and its scalar orbit under $\mathbb F_3^\times$ is free.

This is a quartic local-domain example.  A root of $t^4-27$ has
$3$-adic valuation $3/4$, so its field degree is divisible by four;
the polynomial is therefore irreducible over $\Q_3$, and hence also
over $\Q$.  Its normalization is $O[\pi]$, where
$\pi=3/t$, $\pi^4=3$, and $t=\pi^3$.  The ring $O[\pi]$ is a DVR
with uniformizer $\pi$.
\end{example}

